\documentclass[10pt]{article}
\usepackage[T1]{fontenc}
\usepackage[utf8]{inputenc}
\usepackage[french]{babel}
\usepackage{amsmath,amssymb,mathrsfs,mathpazo}
\usepackage{geometry}
\geometry{a4paper,margin=2cm}
\usepackage{graphicx,array,tabularx,tabularray,colortbl,multirow,multicol,booktabs,enumitem,pifont,ulem,fancybox,diagbox,cancel}
\usepackage{pstricks,pstricks-add,pst-plot,pst-tree,pst-node,pst-text,pst-func,pst-math,pst-bspline,pst-3dplot,pst-eucl,pst-solides3d}
\let\pscalculate\undefined
\usepackage{pst-bezier}
\usepackage{tikz,pgfplots,tkz-tab}
\usetikzlibrary{arrows,arrows.meta,patterns,calc,positioning,trees,decorations.markings,shapes.geometric,shapes.misc,shapes.arrows,intersections,decorations.pathmorphing,decorations.pathreplacing,backgrounds,shadows,babel}
\pgfplotsset{compat=1.18}
\usepgfplotslibrary{fillbetween}
\usepackage[np]{numprint}
\usepackage{listings,makecell,pas-tableur,twoopt}
\let\barre\undefined
\usepackage[pstricks]{tablvar}
\definecolor{darkgreen}{RGB}{0,100,0}
\providecommand{\red}{\color{red}}
\providecommand{\blue}{\color{blue}}
\providecommand{\green}{\color{green}}
\providecommand{\black}{\color{black}}
\providecommand{\white}{\color{white}}
\providecommand{\gray}{\color{gray}}
\providecommand{\cyan}{\color{cyan}}
\providecommand{\magenta}{\color{magenta}}
\providecommand{\yellow}{\color{yellow}}
\providecommand{\violet}{\color{violet}}
\providecommand{\orange}{\color{orange}}
\providecommand{\euro}{\texteuro}
\providecommand{\EUR}{\texteuro}
\providecommand{\ieme}{\textsuperscript{e}}
\providecommand{\degres}{\ensuremath{^\circ}}
\providecommand{\degree}{\ensuremath{^\circ}}
\providecommand{\vv}[1]{\overrightarrow{#1}}
\providecommand{\Coord}[1]{#1}
\providecommand{\decosix}{}
\providecommand{\starredbullet}{\ensuremath{\bullet}}
\pagestyle{empty}
\setlength{\parindent}{0pt}
\usepackage{adjustbox}\DeclareUnicodeCharacter{2236}{:}\DeclareUnicodeCharacter{2190}{\ensuremath{\leftarrow}}
\providecommand{\R}{\mathbb{R}}
\providecommand{\N}{\mathbb{N}}
\providecommand{\Z}{\mathbb{Z}}
\providecommand{\Q}{\mathbb{Q}}
\providecommand{\D}{\mathbb{D}}
\providecommand{\C}{\mathbb{C}}
\providecommand{\e}{{\mathrm{e}}}
\providecommand{\vect}[1]{\overrightarrow{#1}}
\providecommand{\vectt}[1]{\overrightarrow{\mathrm{#1}}}
\providecommand{\vvt}[1]{\overrightarrow{\mathrm{#1}}}
\providecommand{\np}[2][]{{#2~#1}}
\providecommand{\vv}[1]{{\overrightarrow{#1}}}
\renewcommand{\d}{{\mathrm{d}}}
\providecommand{\degres}{{^\circ}}
\providecommand{\degre}{{^\circ}}
\providecommand{\blue}{\color{blue}}
\providecommand{\red}{\color{red}}
\providecommand{\enskip}{\quad}
\providecommand{\llbracket}{\lbrack\!\lbrack}
\providecommand{\rrbracket}{\rbrack\!\rbrack}
\providecommand{\euro}{{\text{€}}}
\providecommand{\no}{{\text{n°}}}
\providecommand{\interval}[2]{{[#1;#2]}}
\providecommand{\intervFO}[2]{{[#1;#2[}}
\providecommand{\dotfill}{\ldots}

\providecommand{\ds}{\displaystyle}
\providecommand{\textding}[1]{#1}
\providecommand{\Oij}{\(\mathrm{O};\vec{\imath},\vec{\jmath}\)}
\providecommand{\Oijk}{\(\mathrm{O};\vec{\imath},\vec{\jmath},\vec{k}\)}
\providecommand{\Ouv}{\(\mathrm{O};\vec{u},\vec{v}\)}
\providecommand{\intervFF}[2]{\ensuremath{[#1;#2]}}
\renewcommand{\euro}{\ifmmode\text{\texteuro}\else\texteuro\fi}
\DeclareUnicodeCharacter{00B0}{\ifmmode^{\circ}\else\textdegree\fi}
\DeclareUnicodeCharacter{20AC}{\ifmmode\text{\texteuro}\else\texteuro\fi}
% Énoncés : Éducation nationale. Transcription et corrigés : APMEP, Denis Vergès et contributeurs. Découpage et classement : CoopMaths, Sébastien Lozano, Éric Elter et contributeurs. Adaptation de la présentation : Maths Explorer.
% CC BY-SA : https://coopmaths.fr/www/about/
% https://coopmaths.fr/alea/static/bac/2024/tex/bac_2024_06_sujet2_metropole_devoile_2_cor.tex
\begin{document}
\noindent\textit{APMEP / CoopMaths. CC BY-SA. Présentation adaptée par Maths Explorer.}\par\medskip
\medskip

\emph{Pour chacune des affirmations suivantes, indiquer si elle est vraie ou fausse. Chaque réponse doit être justifiée. Une réponse non justifiée ne rapporte aucun point.
Les quatre questions de cet exercice sont indépendantes.}

\medskip


\begin{center}A$(0~;~ 4~;~-1)$, \quad B(6~;~1~;~5) \quad et \quad C$(6~;~-2~;~-1)$.\end{center} 



\medskip

\textbf{Affirmation 1 :} Le vecteur $\vect{n}\begin{pmatrix}2\\2\\-1\end{pmatrix}$ est un vecteur normal au plan (ABC).

\medskip

On a $\vect{\text{AB}}\begin{pmatrix}6\\-3\\6\end{pmatrix}$ et $\vect{\text{AC}}\begin{pmatrix}6\\-6\\0\end{pmatrix}$.

Or  $\vect{n} \cdot \vect{\text{AB}} = 12 - 6 - 6 = 0$ et $\vect{n} \cdot \vect{\text{AC}} = 12 - 12  + 0  = 0$.

Or il est admis que les points A, B et C ne sont pas alignés, donc les vecteurs $\vect{\text{AB}}$ et $\vect{\text{AC}}$ ne sont pas colinéaires.

Donc le vecteur $\vect{n}$ est orthogonal à deux vecteurs non colinéaires du plan (ABC) : c'est un vecteur normal à ce plan : affirmation \textbf{Vraie}.
\textbf{Affirmation 2 :} 






On sait que :

$M(x~;~y~;~z) \in (\text{AB}) \iff \vect{\text{A}M} = u\vect{\text{AB}}$.

Avec $\vect{\text{A}M}\begin{pmatrix}x - 0\\y - 4\\z - (-1-\end{pmatrix}$ et $\vect{\text{AB}}\begin{pmatrix}6\\-3\\6\end{pmatrix}$, vecteur qui est colinéaire à $\dfrac13\vect{\text{AB}}\begin{pmatrix}2\\-1\\2\end{pmatrix}$.

Donc $M(x~;~y~;~z) \in (\text{AB}) \iff \left\{\begin{array}{l c l}
x - 0&=& 2u\\
y - 4&=&  - u\\
z - (-1)&=&2u
\end{array}\right. \:\text{où } u \in \R. \iff \left\{\begin{array}{l c l}
x &=& 2u\\
y &=&4  - u\\
z &=&- 1 + 2u
\end{array}\right. \:\text{où } \in \R.$

On pose $u = t + 1$ et on obtient un autre système paramétrique :

$M(x~;~y~;~z) \in (\text{AB}) \iff \left\{\begin{array}{l c l}
x&=&2t + 2\\y &=& 3 - t\\z&=&2t + 1
\end{array}\right. \:\text{où } t \in \R.$

L'affirmation est \textbf{Vraie}.

\medskip

\textbf{Affirmation 3 :} 

Puisque $\mathcal{P}$ est orthogonal à la droite(AB) , le vecteur $\vect{\text{AB}}$ ou plus simplement le vecteur $\dfrac13\vect{\text{AB}}\begin{pmatrix}2\\-1\\2\end{pmatrix}$ est un vecteur normal au plan $\mathcal{P}$.

On sait qu'alors $M(x~;~y~;~z) \in \mathcal{P} \iff 2x - y + 2z + d = 0$ avec $d \in \R$.

Comme C$(6~;~-2~;~-1) \in \mathcal{P} \iff 12  + 2 - 2 + d = 0\iff 12 + d = 0 \iff d = - 12$, on obtient finalement :

$M(x~;~y~;~z) \in \mathcal{P} \iff 2x - y + 2z - 12= 0$ : l'affirmation  est \textbf{Fausse}.

\bigskip



\[\mathcal{D}\quad \left\{\begin{array}{l c l}
x&=&3+t\\
y&=&1+t\\
z&=&2+t
\end{array}\right. \:\text{où } t \in \R \:\:;\quad
\mathcal{D}'\quad \left\{\begin{array}{l c l}
x&= &2t'\\
y&=&4 - t'\\
z&=&-1 + 2t'
\end{array}\right. \:\text{où } t' \in \R.\]

\textbf{Affirmation 4 :} $\mathcal{D}$ et $\mathcal{D}'$ ne sont pas coplanaires.

La droite $\mathcal{D}$ contient le point D(3~;~1~;~2) et a pour vecteur directeur $\vect{d}\begin{pmatrix}1\\1\\1\end{pmatrix}$ et la droite $\mathcal{D}'$ contient le point E(0~;~4~;~-1) et a pour vecteur directeur $\vect{d'}\begin{pmatrix}2\\-1\\2\end{pmatrix}$.

Les vecteurs $\vect{d}$ et $\vect{d}'$ ne sont pas colinéaires donc les droites ne sont pas parallèles.

Si elles ont point commun les coordonnées de celui-ci vérifient les deux systèmes d'où le nouveau système :

$\left\{\begin{array}{l c l}
3+t&=&2t'\\
1+t&=&4 - t'\\
2+t&=& - 1+ 2t'
\end{array}\right.\iff \left\{\begin{array}{l c l}
t&=&2t' - 3\\
t&=&3 - t'\\
t&=& - 3+ 2t'
\end{array}\right. \Rightarrow \left\{\begin{array}{l c l}
2t' - 3&=&3 - t'\\
3 - t'&=& - 3+ 2t'
\end{array}\right. \iff $

$\left\{\begin{array}{l c l}
3t' &=&6 \\
6&=& 3t'
\end{array}\right. \iff \left\{\begin{array}{l c l}
t' &=&2 \\
2&=& t'
\end{array}\right.$ et comme $t = 3 - t' = 3 - 2 = 1$. Le système a une solution.

Conclusion : les droites  sont sécantes au point de coordonnées (4~;~2~;~3) ; elles sont coplanaires : l'affirmation 4 est \textbf{Fausse}
\bigskip

\end{document}
